% ============================================================================
% Fichier  : partie2/P02_C08_ebios.tex
% Titre    : ISO/IEC 27005 et EBIOS Risk Manager
% Auteur   : MINKA MI NGUIDJOÏ Thierry Emmanuel
% Version  : 1.0 -- novembre 2026
% ============================================================================
% Sources : ISO/IEC 27005:2022 ; ISO/IEC 27001:2022 ; ANSSI, EBIOS Risk
% Manager (2018). Cotations de l'écosystème de LiZbiZ : cas fictif.
% ============================================================================

\chapter{ISO/IEC 27005 et EBIOS Risk Manager}
\label{chap:ebios}
\chaptermark{ISO/IEC 27005 et EBIOS RM}

\epigraph{%
    « L'attaquant ne vise pas vos serveurs~;\\
    il vise ce qu'ils permettent. »%
}{Principe du cours}

\niveauChapitre{Intermédiaire}{Chapitres~\ref{chap:cadres}
à~\ref{chap:matrices}~; cours SEC3266, chapitre sur les risques.}

\begin{boxobjectifs}
\begin{itemize}
    \item Situer l'ISO/IEC~27005 dans le système de management de
          l'ISO/IEC~27001.
    \item Distinguer l'approche par les événements et l'approche par les
          actifs.
    \item Conduire les cinq ateliers d'EBIOS Risk Manager.
    \item Coter la menace des parties prenantes de l'écosystème.
    \item Construire un scénario stratégique et un scénario opérationnel.
\end{itemize}
\end{boxobjectifs}

\bigskip

% ============================================================================
\section*{Question directrice : l'attaque de mars pouvait-elle être écrite d'avance~?}
% ============================================================================

En juillet, Mme~MOUAFO conduit la première étude EBIOS Risk Manager de
LiZbiZ. Au troisième atelier, la cartographie de l'écosystème place SAHEL
INFOGÉRANCE seule dans la zone de danger. Au quatrième, le scénario
opérationnel décrit, étape par étape, l'attaque de mars. Ce chapitre
déroule la méthode, deux séances durant.

\medskip
\textbf{Plan du chapitre.}
\begin{enumerate}
    \item L'ISO/IEC 27005 (section~\ref{sec:c08-27005}).
    \item EBIOS Risk Manager~: principes (section~\ref{sec:c08-principes}).
    \item Ateliers 1 et 2~: socle, sources de risque (section~\ref{sec:c08-a12}).
    \item Atelier 3~: l'écosystème (section~\ref{sec:c08-a3}).
    \item Atelier 4~: les scénarios opérationnels (section~\ref{sec:c08-a4}).
    \item Atelier 5~: le traitement (section~\ref{sec:c08-a5}).
\end{enumerate}

\separateur

% ============================================================================
\section{L'ISO/IEC 27005}
\label{sec:c08-27005}
% ============================================================================

L'\norme{ISO/IEC~27005:2022} donne des recommandations pour gérer les
risques liés à la sécurité de l'information, conformément aux exigences
de l'\norme{ISO/IEC~27001}~\cite{iso27005_2022,iso27001_2022}. Elle
reprend le processus de l'ISO~31000 et l'applique au système de
management de la sécurité de l'information.

\begin{tableaucours}{L{3.4cm}L{4.4cm}L{4.4cm}}{%
    \tableauHeader{Approche} &
    \tableauHeader{Point de départ} &
    \tableauHeader{Usage}
}
Par les événements & Les scénarios stratégiques, les sources de risque et
leurs chemins & Vue de direction, risques majeurs, écosystème \\
Par les actifs & Les actifs, leurs menaces et leurs vulnérabilités &
Analyse détaillée d'un périmètre technique \\
\end{tableaucours}
\vspace{-0.8em}
\begin{center}{\small\itshape Tableau~8.1. Les deux approches de l'ISO/IEC 27005:2022}\end{center}

Dans un système de management ISO/IEC~27001, l'appréciation des risques
conduit au choix des mesures de l'annexe~A et à la \termecle{déclaration
d'applicabilité}, qui justifie chaque mesure retenue ou exclue.

% ============================================================================
\section{EBIOS Risk Manager : Principes}
\label{sec:c08-principes}
% ============================================================================

EBIOS Risk Manager, publiée par l'ANSSI en 2018, est une méthode par
scénarios, compatible avec l'ISO~31000 et l'ISO/IEC~27005~\cite{anssi2018ebios}.
Elle s'organise en cinq ateliers.

\begin{tableaucours}{L{3.2cm}L{4.4cm}L{4.6cm}}{%
    \tableauHeader{Atelier} &
    \tableauHeader{Question} &
    \tableauHeader{Livrable}
}
1. Cadrage et socle & Que protège-t-on, contre quel événement redouté~? &
Valeurs métier, biens supports, événements redoutés et gravité, socle de
sécurité \\
2. Sources de risque & Qui pourrait attaquer, pour quoi~? & Couples source
de risque et objectif visé, retenus ou non \\
3. Scénarios stratégiques & Par quel chemin dans l'écosystème~? &
Cartographie de la menace, scénarios stratégiques, mesures sur
l'écosystème \\
4. Scénarios opérationnels & Par quelles étapes techniques~? & Modes
opératoires, vraisemblance \\
5. Traitement & Quelles mesures, quel risque résiduel~? & Stratégie,
plan d'amélioration continue, risques résiduels \\
\end{tableaucours}
\vspace{-0.8em}
\begin{center}{\small\itshape Tableau~8.2. Les cinq ateliers d'EBIOS Risk Manager}\end{center}

La méthode distingue un \termecle{cycle stratégique} (ateliers 1 à 3,
repris à un rythme pluriannuel ou lors d'un changement majeur) et un
\termecle{cycle opérationnel} (ateliers 4 et 5, repris plus souvent, au
fil de l'évolution des menaces et des mesures).

% ============================================================================
\section{Ateliers 1 et 2 : Socle et Sources de Risque}
\label{sec:c08-a12}
% ============================================================================

\begin{tableaucours}{L{3.4cm}L{5.4cm}c}{%
    \tableauHeader{Valeur métier} &
    \tableauHeader{Événement redouté} &
    \tableauHeader{Gravité}
}
Facturation & Facturation arrêtée plus de trois jours & G4 \\
Suivi des expéditions & Suivi indisponible ou faussé plus d'un jour & G3 \\
Données clients & Divulgation de données personnelles & G4 \\
Paiements & Virement frauduleux de plus de 50~millions & G3 \\
\end{tableaucours}
\vspace{-0.8em}
\begin{center}{\small\itshape Tableau~8.3. Atelier 1~: valeurs métier et événements redoutés de LiZbiZ}\end{center}

\begin{tableaucours}{L{3.6cm}L{5.4cm}L{2.6cm}}{%
    \tableauHeader{Source de risque} &
    \tableauHeader{Objectif visé} &
    \tableauHeader{Décision}
}
Groupe criminel & Lucre par rançon et revente de données & Retenu \\
Complice interne & Fraude au virement & Retenu \\
Ancien salarié & Revente de données clients & Retenu \\
Concurrent & Espionnage des tarifs & Non retenu (motivation faible) \\
Militant & Défiguration du portail & Non retenu (gravité faible) \\
\end{tableaucours}
\vspace{-0.8em}
\begin{center}{\small\itshape Tableau~8.4. Atelier 2~: couples source de risque et objectif visé}\end{center}

% ============================================================================
\section{Atelier 3 : l'Écosystème}
\label{sec:c08-a3}
% ============================================================================

L'attaquant passe souvent par un tiers moins protégé. L'atelier~3 cote
chaque partie prenante sur quatre critères, de 1 à 4.

\begin{boxdef}
\termecle{Niveau de menace d'une partie prenante} (EBIOS RM)~:
\[
\text{Menace} = \frac{\text{Dépendance} \times \text{Pénétration}}
{\text{Maturité cyber} \times \text{Confiance}}
\]
La \emph{dépendance} et la \emph{pénétration} (accès au système) exposent~;
la \emph{maturité cyber} et la \emph{confiance} protègent. Les seuils des
zones de danger, de contrôle et de veille sont fixés par l'étude.
\end{boxdef}

\begin{tableaucours}{L{3.6cm}ccccc}{%
    \tableauHeader{Partie prenante} &
    \tableauHeader{Dép.} &
    \tableauHeader{Pén.} &
    \tableauHeader{Mat.} &
    \tableauHeader{Conf.} &
    \tableauHeader{Menace}
}
SAHEL INFOGÉRANCE & 4 & 4 & 2 & 2 & 4,00 \\
Transporteurs sous-traitants & 2 & 2 & 1 & 2 & 2,00 \\
Éditeur de l'application & 3 & 2 & 3 & 3 & 0,67 \\
Clients reliés par échange de données & 2 & 2 & 2 & 3 & 0,67 \\
Opérateur de Garoua & 3 & 1 & 3 & 3 & 0,33 \\
Banque & 2 & 1 & 4 & 4 & 0,13 \\
\end{tableaucours}
\vspace{-0.8em}
\begin{center}{\small\itshape Tableau~8.5. Atelier 3~: menace des parties prenantes de LiZbiZ}\end{center}

\begin{boxscenario}{Opération LIZBIZ. La cartographie de juillet}
\textbf{Seuils retenus par l'étude}~: danger au-dessus de~2,5~; contrôle
de 0,9 à~2,5~; veille en dessous de~0,9.

\medskip
\textbf{Constatation.} SAHEL INFOGÉRANCE (4,00) est seule en zone de
danger~; les transporteurs (2,00) sont en zone de contrôle~; les autres
parties prenantes sont en veille.

\medskip
\textbf{Scénario stratégique retenu.} Un groupe criminel compromet un
technicien du prestataire, utilise son accès distant pour atteindre le
domaine, chiffre les serveurs et exige une rançon (événement redouté~:
facturation arrêtée plus de trois jours, G4).

\medskip
\textbf{Mesures sur l'écosystème.} Second facteur et bastion pour le
prestataire (pénétration et maturité)~; clauses de sécurité et audit
annuel (confiance). Objectif~: ramener la menace sous~2,5.
\end{boxscenario}

% ============================================================================
\section{Atelier 4 : les Scénarios Opérationnels}
\label{sec:c08-a4}
% ============================================================================

L'atelier~4 décrit le mode opératoire, étape par étape, et cote sa
vraisemblance sur une échelle V1 à V4.

\begin{tableaucours}{L{3.0cm}L{6.4cm}L{2.8cm}}{%
    \tableauHeader{Étape} &
    \tableauHeader{Action de l'attaquant} &
    \tableauHeader{Mesure qui l'arrête}
}
Connaître & Identifier un technicien du prestataire et son fournisseur de
messagerie & Sensibilisation \\
Entrer & Hameçonner, puis se connecter à la passerelle avec le mot de
passe & Second facteur \\
Trouver & Explorer le réseau, trouver un serveur non corrigé & Correctifs,
segmentation \\
Exploiter & Obtenir les droits d'administrateur du domaine & Mots de passe
locaux uniques, bastion \\
Agir & Copier, détruire les sauvegardes, chiffrer & Sauvegardes hors ligne,
détection \\
\end{tableaucours}
\vspace{-0.8em}
\begin{center}{\small\itshape Tableau~8.6. Atelier 4~: scénario opérationnel du rançongiciel}\end{center}

\textbf{Vraisemblance.} Avant mesures~: V4 (quasi certain), chaque étape
étant facile. Après second facteur, correctifs et sauvegardes hors ligne~:
V2, l'attaquant devant franchir plusieurs barrières indépendantes.

% ============================================================================
\section{Atelier 5 : le Traitement}
\label{sec:c08-a5}
% ============================================================================

\begin{boxprocedure}
\textbf{Atelier 5 chez LiZbiZ}
\begin{enumerate}
    \item Choisir, pour chaque risque, une option (réduire, transférer,
          éviter, accepter) au regard de l'appétence.
    \item Inscrire les mesures dans un \termecle{plan d'amélioration
          continue de la sécurité} (PACS)~: responsable, échéance, coût,
          indicateur.
    \item Réévaluer gravité et vraisemblance avec les seules mesures en
          place et vérifiées~: c'est le risque résiduel.
    \item Faire accepter formellement les risques résiduels par la
          direction.
    \item Fixer la date de la prochaine revue du cycle opérationnel.
\end{enumerate}
\end{boxprocedure}

\begin{boxobjet}{R-02 · Usurpation de l'accès distant du prestataire}
\textbf{Scénario~:} stratégique et opérationnel des Tableaux~8.5 et~8.6.
\textbf{Cotation~:} G4, V4 avant mesures~; G4, V2 cible.
\textbf{Traitement~:} second facteur (juillet), bastion (septembre),
clauses et audit (contrat renouvelé en octobre). \textbf{Propriétaire~:}
M.~NGUELE. \textbf{Indicateur~:} niveau de menace de SAHEL INFOGÉRANCE
recoté à chaque revue.
\end{boxobjet}

\begin{boxmetier}{Gestionnaire de risques}
\textbf{Faire écrire le scénario par ceux qui le vivraient.} Mme~MOUAFO
fait rédiger le mode opératoire par l'informatique et le prestataire
ensemble~: c'est en décrivant l'attaque étape par étape que chacun voit la
barrière qui manque.
\end{boxmetier}

\separateur

\begin{boxresume}
\textbf{Ce qu'il faut retenir de ce chapitre~:}
\begin{itemize}
    \item \textbf{L'ISO/IEC~27005} applique l'ISO~31000 au système de
          management de la sécurité de l'information~; deux approches, par
          les événements ou par les actifs.
    \item \textbf{EBIOS RM}~: cinq ateliers, un cycle stratégique et un cycle
          opérationnel.
    \item \textbf{L'écosystème} se cote~: dépendance × pénétration, divisé
          par maturité × confiance.
    \item \textbf{Le scénario opérationnel} montre, étape par étape, la
          barrière qui manque.
    \item \textbf{Le traitement} produit un plan, un risque résiduel et une
          acceptation formelle.
\end{itemize}
\end{boxresume}

\begin{boxinfo}
\textbf{Points de contrôle de l'élève.} Je sais~: situer l'ISO/IEC~27005~;
nommer les cinq ateliers~; formuler un couple source de risque et objectif
visé~; calculer le niveau de menace d'une partie prenante~; décrire un
scénario opérationnel et ses barrières.
\end{boxinfo}

\bigskip

\begin{boxexo}[title={\ding{45}\quad Exercice~8.1. Les ateliers \niveaubase}]
Dans quel atelier d'EBIOS RM se rangent~: la liste des valeurs métier~; la
cotation de la maturité du prestataire~; le plan d'amélioration continue~?
\end{boxexo}

\begin{boxexo}[title={\ding{45}\quad Exercice~8.2. Calculer la menace \niveaubase}]
Une société de gardiennage a une dépendance de 2, une pénétration de 1,
une maturité de 1 et une confiance de 2. Calculez sa menace et sa zone.
\end{boxexo}

\begin{boxexo}[title={\ding{45}\quad Exercice~8.3. Les couples \niveaubase}]
Formulez deux couples source de risque et objectif visé pour l'ANPP.
\end{boxexo}

\begin{boxexo}[title={\ding{45}\quad Exercice~8.4. Recoter le prestataire \niveaumoyen}]
Après les mesures, la pénétration de SAHEL INFOGÉRANCE passe à 2, sa
maturité à 3 et la confiance à 3. Recalculez sa menace. Quitte-t-elle la
zone de danger~?
\end{boxexo}

\begin{boxexo}[title={\ding{45}\quad Exercice~8.5. Les transporteurs \niveaumoyen}]
Les transporteurs sont en zone de contrôle. Quelles mesures proposer, sur
quels critères agissent-elles~?
\end{boxexo}

\begin{boxexo}[title={\ding{45}\quad Exercice~8.6. Écarter \niveaumoyen}]
Pourquoi le couple « concurrent, espionnage des tarifs » n'a-t-il pas été
retenu~? Dans quelles circonstances faudrait-il le réexaminer~?
\end{boxexo}

\begin{boxexo}[title={\ding{45}\quad Exercice~8.7. Un scénario opérationnel \niveaumoyen}]
Décrivez le scénario opérationnel « complice interne, fraude au
virement » en cinq étapes, avec une barrière par étape.
\end{boxexo}

\begin{boxexo}[title={\ding{45}\quad Exercice~8.8. Barrières indépendantes \niveaumoyen}]
Pourquoi plusieurs barrières \emph{indépendantes} réduisent-elles la
vraisemblance plus qu'une seule barrière très forte~?
\end{boxexo}

\begin{boxexo}[title={\ding{45}\quad Exercice~8.9. L'atelier 1 de la banque \niveauexpert}]
Conduisez l'atelier~1 pour la BANQUE SANAGA~: quatre valeurs métier, leurs
biens supports, les événements redoutés et leur gravité.
\end{boxexo}

\begin{boxexo}[title={\ding{45}\quad Exercice~8.10. Déclaration d'applicabilité \niveauexpert}]
À partir du Tableau~8.6, indiquez quatre mesures de l'annexe~A de
l'ISO/IEC~27001 que LiZbiZ doit retenir, avec leur justification.
\end{boxexo}

\bigskip

\begin{boxinfo}
\textbf{Pour aller plus loin.}
\begin{itemize}
    \item ANSSI, \emph{EBIOS Risk Manager}~\cite{anssi2018ebios}.
    \item ISO/IEC~27005~\cite{iso27005_2022}~; ISO/IEC~27001~\cite{iso27001_2022}.
\end{itemize}
\end{boxinfo}

\begin{boxinfo}
\textbf{Transition.} Une vraisemblance V2 ou V4 reste une cote. Le
Chapitre~\ref{chap:quantitatif} apprend à estimer le risque en fréquences
et en montants, avec la méthode FAIR.
\end{boxinfo}
